\documentclass[12pt,a4paper]{article}

\usepackage[T1]{fontenc}
\usepackage[utf8]{inputenc}
\usepackage{lmodern}
\usepackage{geometry}
\usepackage{amsmath}
\usepackage{amssymb}
\usepackage{hyperref}
\usepackage{xcolor}
\usepackage{booktabs}
\usepackage{array}
\usepackage{parskip}
\usepackage{microtype}
\usepackage{titlesec}
\usepackage{fancyhdr}
\usepackage{mdframed}
\usepackage{enumitem}
\usepackage{graphicx}   % needed for \includegraphics when replacing placeholders
\usepackage{tikz}       % used to draw the placeholder boxes
\usepackage{caption}
\usepackage{float}

\geometry{margin=2.5cm}

% -----------------------------------------------------------------------
% Colors
% -----------------------------------------------------------------------
\definecolor{sectioncolor}{RGB}{30,60,140}
\definecolor{boxbg}{RGB}{240,245,255}
\definecolor{boxframe}{RGB}{30,60,140}
\definecolor{notebg}{RGB}{255,248,230}
\definecolor{noteframe}{RGB}{180,120,0}
\definecolor{exbg}{RGB}{235,248,235}
\definecolor{exframe}{RGB}{0,120,60}
\definecolor{plbg}{RGB}{248,248,248}       % placeholder background
\definecolor{plframe}{RGB}{160,160,160}    % placeholder border
\definecolor{pltext}{RGB}{100,100,100}     % placeholder text

% -----------------------------------------------------------------------
% Section formatting
% -----------------------------------------------------------------------
\titleformat{\section}{\large\bfseries\color{sectioncolor}}{{\thesection.}}{0.5em}{}[\titlerule]
\titleformat{\subsection}{\normalsize\bfseries\color{sectioncolor}}{{\thesubsection.}}{0.5em}{}
\titleformat{\subsubsection}{\normalsize\itshape\color{sectioncolor}}{}{0em}{}

% -----------------------------------------------------------------------
% Callout boxes
% -----------------------------------------------------------------------
\newmdenv[
  backgroundcolor=boxbg, linecolor=boxframe, linewidth=1.5pt, roundcorner=4pt,
  innertopmargin=8pt, innerbottommargin=8pt, innerleftmargin=12pt, innerrightmargin=12pt,
  skipabove=10pt, skipbelow=10pt
]{keyresult}

\newmdenv[
  backgroundcolor=notebg, linecolor=noteframe, linewidth=1.2pt, roundcorner=4pt,
  innertopmargin=6pt, innerbottommargin=6pt, innerleftmargin=10pt, innerrightmargin=10pt,
  skipabove=8pt, skipbelow=8pt
]{notebox}

\newmdenv[
  backgroundcolor=exbg, linecolor=exframe, linewidth=1.2pt, roundcorner=4pt,
  innertopmargin=6pt, innerbottommargin=6pt, innerleftmargin=10pt, innerrightmargin=10pt,
  skipabove=8pt, skipbelow=8pt
]{exbox}

% -----------------------------------------------------------------------
% GRAPHIC PLACEHOLDER COMMAND
%
% Usage:  \placeholder[width]{height}{label}{description}
%
%   width       (optional) default \linewidth
%   height      TikZ dimension, e.g. 5cm
%   label       short bold tag shown in the top-left corner, e.g. "Fig. 1"
%   description longer italic description of the intended content
%
% To replace with a real figure later, comment out \placeholder and
% uncomment a standard \includegraphics + \caption block.
% -----------------------------------------------------------------------
\newcommand{\placeholder}[4][\linewidth]{%
  \begin{figure}[H]
    \centering
    \begin{tikzpicture}
      \draw[fill=plbg, draw=plframe, line width=1pt, dashed]
            (0,0) rectangle (#1, #3);
      % diagonal cross
      \draw[plframe, line width=0.4pt] (0,0) -- (#1,#3);
      \draw[plframe, line width=0.4pt] (#1,0) -- (0,#3);
      % label tag
      \node[anchor=north west, font=\small\bfseries\color{plframe},
            inner sep=4pt] at (0,#3) {#2};
      % description
      \node[anchor=center, font=\small\itshape\color{pltext},
            text width=0.75\linewidth, align=center]
            at ({0.5*#1},{0.5*#3}) {#4};
    \end{tikzpicture}
  \end{figure}
}

% -----------------------------------------------------------------------
% Header / footer
% -----------------------------------------------------------------------
\pagestyle{fancy}
\fancyhf{}
\rhead{\small Soft X-Ray Grating Monochromators: Theory and Simulation}
\lhead{\small\itshape After Reininger (2016, 2025)}
\cfoot{\thepage}
\renewcommand{\headrulewidth}{0.4pt}

% -----------------------------------------------------------------------
% Macros
% -----------------------------------------------------------------------
\newcommand{\al}{\alpha}
\newcommand{\be}{\beta}
\newcommand{\la}{\lambda}
\newcommand{\RP}{\mathrm{RP}}

% =======================================================================
\begin{document}
% =======================================================================

\begin{titlepage}
  \centering
  \vspace*{2cm}
  {\huge\bfseries Advances in Soft X-Ray Beamline Design:\\[0.4em]
   Grating Monochromator Simulations}\\[1.5em]
  {\large\itshape Theory, Aberration Analysis, and Ray-Tracing Examples}\\[2em]
  {\normalsize Based on lectures by \textbf{Ruben Reininger}\\
   (APS/ANL), originally presented at the\\
   \textit{Software for Optical Simulations Workshop},\\
   Trieste, October 2016;\\
   adapted by M.\ Sanchez del Rio,\\
   DESY Hamburg, May 2025}\\[4em]
  {\normalsize\today}
  \vfill
\end{titlepage}

\tableofcontents
\newpage

% =======================================================================
\section{Introduction and Scope}
% =======================================================================

This document summarises the theoretical framework and ray-tracing simulation examples
presented by Reininger for soft X-ray (SXR) grating monochromators.  Four main instrument
types are treated in order of increasing optical sophistication:

\begin{enumerate}
  \item Spherical Grating Monochromator (SGM)
  \item Vertically collimated Plane Grating Monochromator (vertical PGM)
  \item Follath collimated PGM (horizontal deflection, fixed focal plane)
  \item Focusing Variable Line Spacing PGM (FVLS-PGM)
\end{enumerate}

All simulations use the OASYS/SHADOW ray-tracing environment.  Unless otherwise stated,
the source parameters correspond to the (then) present APS undulator: a 2\,m in-vacuum
device with realistic emittance, divergence, and energy spread.

\placeholder{4cm}{Fig.~1}{%
  Overview schematic of the four monochromator types covered in this document
  (SGM, vertical PGM, Follath PGM, FVLS-PGM), showing their relative complexity
  and key optical components.}

% =======================================================================
\section{Theoretical Foundation: Optical Path Length and Fermat's Principle}
% =======================================================================

\subsection{The Optical Path Length Function}

The foundation for all grating monochromator designs is the \textbf{optical path length
function} $F$ for a toroidal grating, expanded in terms of the sagittal coordinate $w$
(across the groove direction) and the meridional coordinate $\ell$ (along the grooves):

\begin{equation}
  F = F_{00} + F_{01}\,w
    + \tfrac{1}{2}F_{20}\,w^2
    + \tfrac{1}{2}F_{02}\,\ell^2
    + \tfrac{1}{2}F_{30}\,w^3
    + \tfrac{1}{2}F_{21}\,w\,\ell^2
    + \cdots
  \label{eq:OPL}
\end{equation}

\textbf{Fermat's principle} requires that, for each diffracted order, $\delta F = 0$ with
respect to variations in the ray coordinates.  Setting each aberration coefficient to zero
yields the design conditions for a perfect (aberration-free) monochromator.

\placeholder{5cm}{Fig.~2}{%
  Geometry of a toroidal grating showing source point $S$, grating vertex, image
  point $S'$, sagittal coordinate $w$, meridional coordinate $\ell$,
  angles of incidence $\alpha$ and diffraction $\beta$, and radii of curvature
  $R$ (meridional) and $\rho$ (sagittal).}

\subsection{Meaning of the Coefficients}

\begin{center}
\renewcommand{\arraystretch}{1.5}
\begin{tabular}{cll}
\toprule
\textbf{Term} & \textbf{Name} & \textbf{Effect when non-zero} \\
\midrule
$F_{00}$ & Central path length    & Constant (irrelevant for focusing) \\
$F_{01}$ & Grating equation term  & Selects wavelength; always present \\
$F_{20}$ & Defocus (meridional)   & Shifts focus along optical axis \\
$F_{02}$ & Defocus (sagittal)     & Sagittal focus mismatch \\
$F_{30}$ & Primary coma           & Asymmetric blur at exit slit \\
$F_{21}$ & Mixed (coma--sagittal) & Couples sagittal and meridional aberrations \\
\bottomrule
\end{tabular}
\end{center}

In practice, the key design goals are:
\begin{itemize}
  \item $F_{20} = 0$: meridional focus at the exit slit position.
  \item $F_{30} = 0$: coma-free operation (often achievable at one energy only).
\end{itemize}

\subsection{Grating Equation and Resolving Power}

The grating equation in diffraction order $m$ is:
\begin{equation}
  m\la = d\,(\sin\al - \sin\be),
  \label{eq:grating}
\end{equation}
where $d$ is the groove spacing, $\al$ is the angle of incidence (positive above normal),
and $\be$ is the diffraction angle (negative in the SGM convention used here, where
$m = -1$ in SHADOW notation corresponds to the physically relevant first order).

The theoretical resolving power is determined by the exit-slit-limited resolution:
\begin{keyresult}
\begin{equation}
  \RP = \frac{\la}{\Delta\la}
       = \frac{r'\,m\cos\be}{d\cdot\Delta_{\rm slit}},
  \label{eq:RP}
\end{equation}
\end{keyresult}
where $r'$ is the grating-to-slit distance and $\Delta_{\rm slit}$ is the FWHM of the
beam at the slit plane.

\subsection{Image Size: Sanity Check Formula}

A key consistency check in all simulations is the predicted image size at the exit slit,
obtained by geometrical demagnification:

\begin{equation}
  \Delta_z = 2.35\,\sigma_z\,\frac{\cos\al}{\cos\be}\,\frac{r'}{r}
  \quad\text{(vertical/meridional plane)},
  \label{eq:sanity_z}
\end{equation}
\begin{equation}
  \Delta_x = 2.35\,\sigma_x\,\frac{r'}{r}
  \quad\text{(horizontal/sagittal plane)},
  \label{eq:sanity_x}
\end{equation}
where $\sigma_z$, $\sigma_x$ are the source RMS sizes and $r$, $r'$ are the
source-to-grating and grating-to-image distances.  Agreement between the ray-tracing
result and Eqs.~\eqref{eq:sanity_z}--\eqref{eq:sanity_x} confirms that the optical system
is properly focused.

\placeholder{3.5cm}{Fig.~3}{%
  Diagram illustrating the demagnification geometry: source of size $\sigma_z$
  at distance $r$, grating with angles $\alpha$ and $\beta$, and image of size
  $\Delta_z$ at distance $r'$, showing the $(\cos\alpha/\cos\beta)$ asymmetry factor.}

% =======================================================================
\section{Spherical Grating Monochromator (SGM)}
% =======================================================================

\subsection{Optical Layout}

The SGM is the simplest SXR monochromator: a single \textbf{concave spherical grating}
works simultaneously as the dispersing and focusing element.  The geometry is defined by:

\begin{itemize}
  \item Source-to-grating distance $r$.
  \item Grating-to-slit distance $r'$ (movable mechanically).
  \item Bragg-like incidence angle $\al$ and diffraction angle $\be$, both measured from
        the grating normal.
  \item Diffraction order $m = -1$.
\end{itemize}

The defocus condition $F_{20} = 0$ gives the \textbf{grating equation for focusing}:

\begin{equation}
  \frac{\cos^2\al}{r} + \frac{\cos^2\be}{r'} = \frac{\cos\al + \cos\be}{R},
  \label{eq:SGM_focus}
\end{equation}
where $R$ is the grating radius of curvature.  At each energy, $\al$ and $\be$ change
(coupled through Eq.~\eqref{eq:grating}), so the focal distance $r'$ also changes.
Therefore, \textbf{the SGM requires a movable exit slit} to track the focus.

\placeholder{4.5cm}{Fig.~4}{%
  SGM optical layout: source at distance $r$, single spherical grating (radius $R$),
  and movable exit slit at distance $r'(E)$.  Indicate the angles $\alpha$ and $\beta$,
  and the energy-dependent slit travel range.}

\begin{notebox}
\textbf{SGM design rule:} Because $r'(E)$ is energy-dependent, a fixed exit slit position
can only be optimal at one energy.  At all other energies, coma and defocus degrade the
resolution.  The exit slit must be translated to follow the focus.
\end{notebox}

\subsection{Ray-Tracing Example at 1000\,eV}

\begin{exbox}
\textbf{File:} \texttt{01\_SGM\_1000eV.ows}\\[4pt]
\textbf{Parameters:} $r = 30$\,m, $r' \approx 9.9$\,m, $\al = 87.3°$, $\be = 89.1°$,
diffraction order $m = -1$.

\medskip
\textbf{Results:}
\begin{itemize}
  \item Predicted vertical image size (sanity check):
        $\Delta_z = 2.35\,\sigma_z\,(\cos\al/\cos\be)\,(r'/r) \approx 36\,\mu$m.
  \item Predicted horizontal image size: $\Delta_x \approx 217\,\mu$m.
  \item Achieved resolving power: $\RP \approx 18\,000$ at 1000\,eV with 55\,meV FWHM
        energy resolution.
  \item Ray-tracing confirms both sanity check values, validating the optical layout.
\end{itemize}

\medskip
\textbf{Slit scan:} Resolution is \emph{source limited} for slit openings
$\lesssim 20\,\mu$m; for larger openings ($\gtrsim 50\,\mu$m) the resolution is
\emph{slit limited}, degrading linearly with slit width.
\end{exbox}

\placeholder{5cm}{Fig.~5}{%
  SHADOW ray-tracing results for SGM at 1000\,eV.  Left: footprint and spot diagram at
  the exit slit plane ($\Delta_z \approx 36\,\mu$m, $\Delta_x \approx 217\,\mu$m).
  Right: resolving power vs.\ slit opening, showing source-limited and slit-limited
  regimes with the transition near $20\,\mu$m.}

\subsection{Ray-Tracing Example at 1600\,eV}

\begin{exbox}
\textbf{File:} \texttt{02\_SGM\_1600eV.ows}\\[4pt]
\textbf{Parameters:} $\al = 87.63°$, $\be = -88.77°$.  Source: $S_x = 280\,\mu$m,
$S_z = 154\,\mu$m, $S_x' = 18\,\mu$rad, $S_z' = 14\,\mu$rad.

\medskip
\textbf{Key finding --- focus shift with energy:}\\
The sanity-check formula predicts $\Delta_z \approx 20\,\mu$m at the nominal slit
position, but the ray-tracing footprint shows the beam is \emph{not at focus} at 1600\,eV
if the exit slit remains at the 1000\,eV position.  Moving the exit slit to
$r' = 11\,046$\,mm (satisfying $F_{20}(1600\,\mathrm{eV}) = 0$) restores focus and
gives $\Delta_z \approx 22\,\mu$m, consistent with the sanity check.

\medskip
\textbf{Conclusion:} This example confirms that the SGM requires slit translation with
energy.  The exit slit movement \emph{corrects defocus but does not correct coma}.
\end{exbox}

\placeholder{5cm}{Fig.~6}{%
  SGM at 1600\,eV: comparison of slit-plane footprints (a)~at the fixed 1000\,eV slit
  position (defocused, asymmetric) and (b)~after slit translation to $r' = 11\,046$\,mm
  (in focus, $\Delta_z \approx 22\,\mu$m).  The residual coma asymmetry remains in
  both cases.}

% =======================================================================
\section{Plane Grating Monochromator (PGM)}
% =======================================================================

\subsection{Key Properties of a Plane Grating}

A plane grating (without variable line spacing) \textbf{does not focus}.  Its role is
purely dispersive.  The focusing is delegated entirely to curved mirrors upstream and
downstream of the grating.  The important consequence is a \textbf{fixed focal plane}:
the image distance is independent of the grating rotation.

The grating magnification factor is defined as:
\begin{equation}
  c = \frac{\cos\be}{\cos\al},
  \label{eq:c_factor}
\end{equation}
and the grating equation becomes:
\begin{equation}
  1 - \frac{m\la}{d\sin\be} = \frac{1}{c^2}.
  \label{eq:PGM_grating}
\end{equation}

The parameter $c$ controls the angular compression/expansion between the collimated beam
and the diffracted beam, and thus the \textbf{grating magnification} $c^{-1}$ can be
tuned by rotating the plane mirror that delivers the collimated beam.

\placeholder{4cm}{Fig.~7}{%
  Schematic of the plane grating geometry.  Show the plane mirror (angle $\theta_M$)
  delivering the collimated beam at incidence angle $\alpha$ to the plane grating, and
  the diffracted beam at angle $\beta$ being focused by the downstream cylindrical mirror.
  Label the magnification factor $c = \cos\beta/\cos\alpha$.}

\subsection{Collimated PGM: Vertical Deflection Layout}

\begin{exbox}
\textbf{File:} \texttt{03\_CollimatedPGM1000eVc2\_Vertical.ows}\\[4pt]
\textbf{Layout} (distances from source): Toroidal mirror at 9\,m $\to$ plane mirror +
plane grating at 28\,m $\to$ cylindrical mirror at 30\,m $\to$ exit slit at 31\,m.

\medskip
\textbf{Optical roles:}
\begin{itemize}
  \item \textbf{Toroidal mirror:} collimates the beam in the meridional (vertical) plane;
        focuses sagittally (horizontally) at the exit slit.
  \item \textbf{Plane mirror + plane grating:} deflect in the vertical plane; the plane
        mirror rotation controls $c$ and hence the wavelength.
  \item \textbf{Cylindrical mirror:} focuses the collimated, dispersed beam onto the
        \emph{fixed} exit slit.
\end{itemize}

\medskip
\textbf{Parameters at 1000\,eV, $c = 2$:} $\al = 88.5°$, $\be = 87.1°$; plane mirror at
$87.8°$ incidence.

\medskip
\textbf{Aberrations:}
\begin{itemize}
  \item \textbf{Coma} is present in the vertical plane because the collimated beam does
        not enter the grating exactly at normal incidence (residual $F_{30}$).
  \item Image size at slit: $\sim 5.8\,\mu$m without slope errors.
  \item \textbf{Slope errors} on the toroidal mirror broaden the slit image significantly.
        The simulation shows that slope errors on the upstream toroid are the dominant
        resolution-limiting factor in this configuration.
\end{itemize}
\end{exbox}

\placeholder{5.5cm}{Fig.~8}{%
  Vertical collimated PGM layout ($c = 2$, 1000\,eV).  Top: beamline schematic with
  distances (toroid at 9\,m, plane mirror + grating at 28\,m, cylinder at 30\,m, slit
  at 31\,m).  Bottom: ray-tracing spot at the exit slit showing coma aberration
  ($\sim 5.8\,\mu$m ideal) and the broadened profile with slope errors on the toroid.}

\placeholder{4cm}{Fig.~9}{%
  Effect of slope errors on the toroidal mirror in the vertical PGM.
  Comparison of slit images: (a)~without slope errors (ideal, $\sim 5.8\,\mu$m) and
  (b)~with realistic r.m.s.\ slope errors, demonstrating the dominant contribution of
  the toroid to the resolution budget.}

\subsection{Follath Collimated PGM: Horizontal Deflection}

The \textbf{Follath PGM} (R.\ Follath and F.\ Senf, NIM A\,390, 388, 1997) uses
\textbf{horizontal} deflection by the plane mirror and grating.  The key advantage is
that, in the horizontal plane, slope errors on the toroidal mirror are much less critical
because the sagittal focusing is less sensitive to meridional slope errors.

\begin{exbox}
\textbf{File:} \texttt{04\_CollimatedPGM1000eVc2.ows}\\[4pt]
\textbf{Same layout as above, but mirror--grating system rotates in the horizontal plane.}

\medskip
\textbf{Key results at $c = 2$, 1000\,eV:}
\begin{itemize}
  \item \textbf{No coma:} the horizontal geometry naturally satisfies $F_{30} \approx 0$
        due to the relative rotation between the plane mirror and the grating in the
        $xz$-plane.
  \item Image size: $\sim 5.8\,\mu$m (same as vertical case without aberrations, but now
        robust against mirror slope errors).
  \item The \textbf{``forgiveness factor''} $\times\theta$: slope errors on the toroidal
        mirror have a much smaller effect in the horizontal geometry.
  \item Simulation with and without slope errors confirms that the Follath layout is
        significantly more tolerant to mirror figure errors.
\end{itemize}

\medskip
\textbf{Effect of increasing $c$:} At $c = 5$ (file
\texttt{05\_CollimatedPGM1000eVc5.ows}), $\al = 89.5°$, $\be = 87.4°$.  The
\textbf{astigmatic coma} grows, scaling as $\sim 5.8/(c/2) = 2.32\,\mu$m for ideal
optics.  With slope errors, the image quality degrades further.
\end{exbox}

\placeholder{5.5cm}{Fig.~10}{%
  Follath horizontal PGM ($c = 2$, 1000\,eV).  Left: beamline plan view showing
  horizontal deflection by the plane mirror and grating.  Centre and right: SHADOW
  footprints at the exit slit without and with toroid slope errors, demonstrating the
  ``forgiveness factor'' --- the slope errors have negligible effect in this geometry.}

\placeholder{4cm}{Fig.~11}{%
  Astigmatic coma growth with increasing magnification factor $c$ in the Follath PGM.
  Slit images at $c = 2$ and $c = 5$ (file \texttt{05\_CollimatedPGM1000eVc5.ows}),
  showing the $5.8/(c/2)$ scaling of the coma contribution.}

% =======================================================================
\section{Focusing Variable Line Spacing PGM (FVLS-PGM)}
% =======================================================================

\subsection{Principle}

The \textbf{FVLS-PGM} (R.\ Reininger and A.\ R.\ B.\ de Castro, NIM A\,538, 760, 2005)
addresses the two main limitations of the collimated PGM:

\begin{enumerate}
  \item \textbf{Heat-load-induced deformations} of the plane mirror degrade the resolution
        by introducing a virtual source shift.
  \item \textbf{Coma} cannot be simultaneously cancelled at all energies with a fixed
        line-spacing grating.
\end{enumerate}

The solution is to use a \textbf{variable line spacing} (VLS) grating, where the groove
density varies along the grating surface as:
\begin{equation}
  \sigma(w) = \sigma_0 + b_2\,w + b_3\,w^2 + \cdots,
  \label{eq:VLS}
\end{equation}
with $\sigma_0 = 1/d_0$ the central line density.  The coefficients $b_2$ and $b_3$
are chosen to:

\begin{itemize}
  \item $b_2$: solve the defocus equation $F_{20} = 0$ at all wavelengths simultaneously,
        by illuminating the grating at the correct angle of incidence (note: in the FVLS
        design, $c$ is \emph{not} a free parameter --- it is determined by $b_2$).
  \item $b_3$: cancel coma $F_{30} = 0$ at one selected wavelength; coma is then
        acceptably small over the full energy range.
\end{itemize}

\begin{keyresult}
\textbf{FVLS design parameters} (example at central energy 600\,eV, $c = 2$):
\begin{align}
  b_2 &= 1.44 \times 10^{-4}\,\mathrm{mm}^{-1}, \notag\\
  b_3 &= 1.3  \times 10^{-8}\,\mathrm{mm}^{-2}. \notag
\end{align}
Coma is zeroed at 600\,eV; solving coma at one wavelength is sufficient because the
residual coma at other energies remains small.
\end{keyresult}

\placeholder{4cm}{Fig.~12}{%
  VLS grating schematic showing the variable groove spacing $d(w) = 1/\sigma(w)$ as a
  function of position $w$ along the grating surface.  Indicate the central density
  $\sigma_0$ and the quadratic variation controlled by $b_2$ and $b_3$.}

\subsection{Optical Layout}

\begin{exbox}
\textbf{File:} \texttt{06\_FVLSPGM.ows}\\[4pt]
\textbf{Distances:} Source $\to$ heat absorber $\to$ plane mirror (28\,m) $\to$ VLS
grating (28.3\,m) $\to$ slit plane (38.3\,m) $\to$ refocusing mirror $\to$ sample
(48.3\,m).\\
Parameters at 1000\,eV: $\al = 88.5°$, $\be = 87.1°$ (same $c = 2$ condition).

\medskip
\textbf{Key feature:} There is no upstream collimating mirror.  The VLS grating itself
provides the meridional focusing at the exit slit, and the horizontal focus is provided
after the slit by a separate refocusing mirror.
\end{exbox}

\placeholder{4.5cm}{Fig.~13}{%
  FVLS-PGM beamline layout.  Show the heat absorber, plane mirror (28\,m), VLS grating
  (28.3\,m), exit slit (38.3\,m), and post-slit refocusing mirror.  Annotate the
  distances and the absence of an upstream collimating mirror compared to the Follath
  PGM.}

\subsection{Heat Load and Correction}

The crucial advantage of the FVLS-PGM emerges when heat-load deformations of the plane
mirror are considered:

\begin{exbox}
\textbf{Heat load simulation:} The plane mirror absorbs $\sim 35$\,W of total power.
The resulting thermo-mechanical deformation produces a maximum surface height deviation of
$\sim 0.40\,\mu$m and slope errors up to $\sim 15\,\mu$rad peak.  The effective radius of
curvature induced is $R \approx -30\,\mathrm{km}$ (very weakly concave).  This shifts the
virtual source position to $\sim 28.52$\,m.

\medskip
\textbf{Without correction:} the deformed mirror creates a defocused, asymmetric image at
the slit --- resolving power degrades substantially.

\medskip
\textbf{Correction strategy:} Because the FVLS grating corrects defocus at all energies
by adjusting the angle of incidence, a \emph{small angular correction} to $\al$ and $\be$
is sufficient to restore the focus:
\[
  \al:\; 88.52843° \to 88.53789°,\quad
  \be:\; 87.05389° \to 87.05860°.
\]
This is equivalent to accounting for the shifted virtual source position in the VLS design
equations.

\medskip
\textbf{Result:} After angular correction, the ray-tracing shows full recovery of the
nominal slit image --- the FVLS-PGM \textbf{self-corrects for heat-load deformations}
with only a minor adjustment of the mirror--grating angles.

\medskip
\textbf{Reference:} R.\ Reininger, NIM A\,649, 139 (2010).
\end{exbox}

\placeholder{5cm}{Fig.~14}{%
  Heat load deformation of the plane mirror (35\,W absorbed power).
  Left: surface height deviation (peak $\sim 0.40\,\mu$m) and slope error profile
  (peak $\sim 15\,\mu$rad) vs.\ mirror length.  Right: induced virtual-source shift
  to $\sim 28.52$\,m from the nominal 28\,m position.}

\placeholder{5.5cm}{Fig.~15}{%
  FVLS-PGM ray-tracing at the exit slit under three conditions:
  (a)~no heat load (reference); (b)~with heat load, uncorrected (defocused, RP degraded);
  (c)~with heat load, corrected by angular adjustment
  ($\Delta\alpha = +0.00946°$, $\Delta\beta = +0.00471°$) --- full RP recovery.}

\begin{notebox}
\textbf{Limitation of the FVLS-PGM:} Because $c$ is fixed by the VLS coefficients
($b_2$), the grating magnification \emph{cannot} be changed during operation.  The Follath
collimated PGM retains the freedom to vary $c$ (and hence the resolving power vs.\ flux
trade-off), but at the cost of reduced heat-load tolerance.
\end{notebox}

% =======================================================================
\section{Comparison of Monochromator Types}
% =======================================================================

\begin{center}
\renewcommand{\arraystretch}{1.7}
\begin{tabular}{p{3.2cm} p{2.5cm} p{2.5cm} p{2.5cm} p{2.5cm}}
\toprule
\textbf{Property}
  & \textbf{SGM}
  & \textbf{Vertical PGM}
  & \textbf{Follath PGM}
  & \textbf{FVLS-PGM} \\
\midrule
Exit slit position
  & Movable (tracks focus)
  & Fixed
  & Fixed
  & Fixed \\
Coma
  & Present
  & Present
  & Eliminated (hor.\ geom.)
  & Zeroed at one $\lambda$ \\
Slope error sensitivity
  & Moderate
  & High (toroid)
  & Low (forgiveness factor)
  & Low \\
Heat load correction
  & Slit movement
  & Slit movement
  & Slit movement
  & Angle correction \\
Variable magnification $c$
  & N/A
  & Yes
  & Yes
  & No (fixed by VLS) \\
Focusing element
  & Grating only
  & Toroid + cylinder
  & Toroid + cylinder
  & VLS grating \\
\bottomrule
\end{tabular}
\end{center}

\placeholder{5cm}{Fig.~16}{%
  Side-by-side comparison of the four monochromator layouts drawn to the same scale,
  highlighting the presence/absence of the upstream collimating mirror, direction of
  deflection (vertical vs.\ horizontal), and the position of the exit slit relative to
  the grating.}

% =======================================================================
\section{Summary and Conclusions}
% =======================================================================

The lecture by Reininger establishes a unified theoretical framework --- the optical path
length expansion of Eq.~\eqref{eq:OPL} and Fermat's principle --- that underpins all
grating monochromator designs for soft X-rays.  The main conclusions are:

\begin{enumerate}[leftmargin=*, label=\arabic*.]
  \item \textbf{The SGM} is conceptually simple but operationally demanding: the exit slit
        must track the focus with energy, and coma cannot be corrected independently.

  \item \textbf{The collimated PGM} decouples focusing from dispersion.  The
        \emph{vertical} layout is coma-affected and sensitive to toroidal mirror slope
        errors.  The \emph{Follath horizontal} layout eliminates coma and is far more
        tolerant to slope errors (the ``forgiveness factor'').

  \item \textbf{The FVLS-PGM} is the most robust design: VLS coefficients simultaneously
        cancel defocus at all energies and coma at the design energy, and the grating
        angles can be slightly adjusted to compensate heat-load deformations of the plane
        mirror without moving the exit slit.

  \item \textbf{OASYS/SHADOW} provides the numerical environment to implement all these
        designs and verify them against the analytical sanity checks (image size,
        resolving power).  The formulas for all monochromator types discussed here are
        included in the OASYS widget library.
\end{enumerate}

\placeholder{4.5cm}{Fig.~17}{%
  Summary figure: resolving power $\RP$ vs.\ photon energy for the four monochromator
  types simulated in this work, plotted on a common axis to allow direct performance
  comparison.  Each curve corresponds to the optimal operating condition of the
  respective instrument.}

% =======================================================================
\newpage
\appendix
\section*{Glossary of Symbols}
\addcontentsline{toc}{section}{Glossary of Symbols}
% =======================================================================

\begin{tabular}{lp{10cm}}
\toprule
\textbf{Symbol} & \textbf{Definition} \\
\midrule
$F$                    & Optical path length function \\
$F_{ij}$               & Coefficients of the OPL expansion ($i$: meridional order,
                          $j$: sagittal order) \\
$w$                    & Coordinate along the groove direction (sagittal) \\
$\ell$                 & Coordinate across the grooves (meridional) \\
$\al$                  & Angle of incidence on the grating (from grating normal) \\
$\be$                  & Angle of diffraction \\
$m$                    & Diffraction order \\
$d$                    & Groove spacing (at grating centre) \\
$\la$                  & Wavelength \\
$r$                    & Source-to-grating distance \\
$r'$                   & Grating-to-image (exit slit) distance \\
$R$                    & Grating radius of curvature (SGM) \\
$c$                    & Magnification factor ($= \cos\be/\cos\al$) \\
$\RP$                  & Resolving power ($= \la/\Delta\la$) \\
$S_x, S_z$             & Source size (horizontal, vertical) RMS \\
$S_x', S_z'$           & Source divergence (horizontal, vertical) RMS \\
$b_2, b_3$             & VLS grating coefficients (line density gradient, curvature) \\
$\sigma_0$             & Central groove density of VLS grating \\
\bottomrule
\end{tabular}

\bigskip
% =======================================================================
\section*{References}
\addcontentsline{toc}{section}{References}
% =======================================================================

\begin{enumerate}
  \item R.\ Follath and F.\ Senf,
        \textit{New plane grating monochromators for third generation synchrotron
        radiation light sources},
        Nucl.\ Instrum.\ Methods Phys.\ Res.\ A \textbf{390}, 388 (1997).
        \href{https://doi.org/10.1016/S0168-9002(97)00401-4}{doi:10.1016/S0168-9002(97)00401-4}

  \item R.\ Reininger and A.\ R.\ B.\ de Castro,
        \textit{The Focusing VLS-PGM: A new monochromator for soft X-rays},
        Nucl.\ Instrum.\ Methods Phys.\ Res.\ A \textbf{538}, 760 (2005).
        \href{https://doi.org/10.1016/j.nima.2004.09.007}{doi:10.1016/j.nima.2004.09.007}

  \item R.\ Reininger,
        \textit{Correction of heat-load-induced deformations in soft X-ray plane grating
        monochromators},
        Nucl.\ Instrum.\ Methods Phys.\ Res.\ A \textbf{649}, 139 (2010).
\end{enumerate}

% =======================================================================
\newpage
\section*{How to Replace Placeholders with Real Figures}
\addcontentsline{toc}{section}{How to Replace Placeholders}
% =======================================================================

Each placeholder in this document is produced by the command:

\begin{verbatim}
\placeholder[width]{height}{Fig. N}{Description of intended content}
\end{verbatim}

To replace a placeholder with a real figure, comment out the \verb|\placeholder| line
and insert a standard \LaTeX\ figure block in its place:

\begin{verbatim}
% \placeholder{5cm}{Fig.~5}{...description...}
\begin{figure}[H]
  \centering
  \includegraphics[width=0.85\linewidth]{fig05_SGM_1000eV_raytracing.pdf}
  \caption{SHADOW ray-tracing results for SGM at 1000\,eV. ...}
  \label{fig:SGM_1000eV}
\end{figure}
\end{verbatim}

Recommended image formats (in order of preference): \texttt{.pdf}, \texttt{.eps},
\texttt{.png}.  Place all figure files in the same directory as the \texttt{.tex} source,
or adjust the path in \verb|\includegraphics| accordingly.

\end{document}